\documentclass[11pt,a4paper]{article}
\usepackage{fontspec}
\setmainfont{Noto Sans CJK SC}[AutoFakeSlant=0.15]
\setsansfont{Noto Sans CJK SC}
\setmonofont{DejaVu Sans Mono}[Scale=0.79]
\IfFontExistsTF{Arial}{\newfontfamily\latinfont{Arial}}{\newfontfamily\latinfont{Liberation Sans}}
\XeTeXlinebreaklocale "zh"
\XeTeXlinebreakskip=0pt plus 1pt
\usepackage{amsmath,amssymb,unicode-math}
\setmathfont{Latin Modern Math}
\usepackage[margin=22mm,headheight=16pt]{geometry}
\usepackage{booktabs,longtable,array,tabularx,graphicx,xcolor,enumitem,listings,fancyhdr,hyperref,multicol}
\definecolor{navy}{HTML}{263B63}
\definecolor{blue}{HTML}{4477AA}
\definecolor{orange}{HTML}{D55E00}
\definecolor{soft}{HTML}{F2F5F9}
\hypersetup{colorlinks=true,linkcolor=navy,urlcolor=blue,citecolor=navy,
 pdftitle={PoreMC 0.6: Cached Host Fields and Exact Delayed CBMC},
 pdfauthor={PoreMC research implementation}}
\setlength{\parindent}{1.7em}
\setlength{\parskip}{0.35em}
\setlength{\emergencystretch}{3em}
\linespread{1.06}
\setlist{itemsep=3pt,topsep=5pt,leftmargin=1.8em}
\renewcommand{\contentsname}{目录}
\renewcommand{\tablename}{表}
\renewcommand{\figurename}{图}
\renewcommand{\refname}{参考文献}
\pagestyle{fancy}\fancyhf{}
\fancyhead[L]{\small\color{navy}PoreMC 0.6}
\fancyhead[R]{\small 固定框架场、参考势生长与精确校正}
\fancyfoot[C]{\small\thepage}
\lstset{basicstyle=\ttfamily\small,breaklines=true,columns=fullflexible,keepspaces=true,
 showstringspaces=false,frame=none,backgroundcolor=\color{soft},
 xleftmargin=6pt,xrightmargin=6pt,aboveskip=7pt,belowskip=7pt}
\newcommand{\code}[1]{\texttt{\small #1}}
\newcommand{\TotalTests}{335}
\newcommand{\RigidSpeed}{16.71}
\newcommand{\FlexibleSpeed}{46.22}

\begin{document}
\begin{titlepage}
\vspace*{18mm}
{\latinfont\fontsize{44}{50}\selectfont\color{navy}\bfseries PoreMC 0.6\par}
\vspace{12mm}
{\fontsize{25}{35}\selectfont 固定框架场与高效巨正则采样\par}
\vspace{5mm}
{\Large 缓存、参考势 CBMC、Boltzmann 插值与精确 Ewald 校正\par}
\vspace{12mm}
\noindent\color{navy}\rule{\textwidth}{1.1pt}\color{black}
\vspace{8mm}
\noindent 从用户给出的 4,416 原子框架与 12 位点客体瓶颈出发，把固定环境的重复计算移出采样循环，把完整倒空间评价移到已选中的候选上。理论推导、原生实现、反例分析和独立测量共同确定最终算法。\par
\vspace{10mm}
\begin{tabular}{@{}p{31mm}p{118mm}@{}}\toprule
固定环境 & LJ 类型场与单位电荷静电场；一次 FFT 建表；二进制缓存跨温度和压力复用。\\[5pt]
采样流程 & 网格参考势构造 CBMC 权重；两阶段接受率恢复原有限 Ewald 目标。\\[5pt]
搜索空间 & 陡峭区域 Boltzmann 权重插值；可选全支持 alias 根位点提议。\\[5pt]
倒空间 & 精确共轭半谱、平移相位更新、接受态与逐分子贡献复用。\\[5pt]
本轮证据 & 单线程完整循环、接受移动吞吐、长链载量与相关性、独立数值核验。\\\bottomrule
\end{tabular}
\vfill
\noindent\textbf{版本}\quad 0.6.0\qquad\textbf{日期}\quad 2026 年 11 月 18 日\par
\noindent\textbf{单位}\quad 能量 $u=U/k_{\mathrm B}$：K；坐标：\AA；电荷：$e$。\par
\end{titlepage}

\section*{实现结果与选择依据}
\addcontentsline{toc}{section}{实现结果与选择依据}
\code{FastGCMC} 新增原生 \code{sampling\_mode='delayed'}。候选阶段通过固定框架场和客体短程作用选择构型；只有通过第一阶段的候选才评价精确长程残差。接受后复用已计算的分子贡献。默认保留原矩形 $k_{\max}$ 集合，网格插值仅定义参考势，最终目标 Hamiltonian 保持不变。

本轮同机、同物理输入、单线程完整循环中，12 位点刚性合成算例的 Boltzmann $32^3$ 路径相对原始 0.5 的配对吞吐提升中位数为 \RigidSpeed{} 倍；12 位点柔性链为 \FlexibleSpeed{} 倍。两类均采用 4,416 框架原子、10 个候选、10 个初始分子和 $k_{\max}=(13,13,13)$。接受移动吞吐、载量和混合检验与速度同时报告。

初始的粗网格能量三线性方案虽然计算很快，却使近壁吸附区域难以进入。程序随后改为陡峭区域的 Boltzmann 权重插值，并用独立分布与长链对比检验其效果。该反例说明，网格应服务于更有效的采样，单独增加尝试数并不足以确定最佳参数。

最终构建的 \TotalTests{} 项测试通过。完整源代码包含实际运行的脚本和原始证据，报告表格直接由这些 JSON 和测试日志生成。用户提供的 RASPA2 HPC 数据作为问题背景保留；本轮新速度是 PoreMC 0.5/0.6 的本机配对测量。

\setcounter{tocdepth}{1}
\begingroup\small\begin{multicols}{2}\tableofcontents\end{multicols}\endgroup
\clearpage

\section{工作负载与加速目标}
用户实测体系包含 4,416 个框架原子、12 位点 p-xylene、10 个 CBMC trial，以及插入、删除、平移、旋转和再生长。0.5 direct Ewald 使用 19,682 个完整矩形倒空间模式；用户记录的 RASPA2 使用约 5,685 个有效模式并复用框架网格。输入中的有效模式定义和原始构型未随本轮消息提供，因此本版分别处理模式代数压缩与原物理模型下的采样流程优化。

\begin{center}\small
\begin{tabular}{@{}lrr@{}}\toprule
用户 HPC 数据 & 原始吞吐 & 150,000 cycles 的给定折算\\\midrule
RASPA2，平均约 5.5 分子/晶胞 & 123.7 cycles/s & 20.2 min\\
RASPA2，平均约 17.6 分子/晶胞 & 75.1 cycles/s & 33.3 min\\
PoreMC 0.5 direct & 89 attempts/s & 9.36 h\\
PoreMC 0.5 PME $32^3$ & 151.5 attempts/s & 5.50 h\\\bottomrule
\end{tabular}\end{center}

上述折算按每循环 20 次尝试，即 $3\times10^6$ 次移动。若循环定义为 $\max(20,N)$，精确工作量应从粒子数轨迹累计，不能只把平均 $N$ 代入。用户表中的两个 RASPA2 点对应约 1,502--2,474 attempts/s 的固定 20 次换算参考。

目标瓶颈可写为 $B\times A\times K$：候选数、客体位点数和倒空间模式数相乘。固定框架场已具有时间不变性，CBMC 又只会保留少数候选。因此本版同时减少每个候选的代价、需要精确评价的候选数，以及精确评价内部的重复项。框架网格与近似接受后精确校正的理论基础分别见文献 \cite{ran2025,christen2005,dubbeldam2013}。

% PoreMC 0.6 theory chapter. Included by the report main file.
% Required: amsmath, amssymb, booktabs, longtable, array, listings; \code{}.

\section{把固定框架、参考提议与目标能量分开组织}
\label{sec:06-decomposition}

固定框架体系具有两个可利用的结构：框架场在整条轨迹上不变，CBMC 的多个候选最终只有一个可能成为新构型。0.6 分别沿这两个方向重构计算：把框架作用变成可复用的空间场；用低成本参考势完成候选选择，再对通过第一阶段的完整分子计算目标势与参考势的残差。网格插值、分拆长程能量的 CBMC 和 delayed acceptance 均有既有理论基础\cite{dubbeldam2013,ran2025,christen2005}。本版的工作落在原生状态、缓存粒度、候选执行顺序和可复核的数值实现。

令 $X_N$ 表示 $N$ 个客体分子的构型，能量变量为 $u=U/k_{\mathrm B}$，单位 K。内部几何使用归一化参考测度 $\mathrm d\nu$，活度 $z$ 单位为 \AA$^{-3}$。目标巨正则测度为
\begin{equation}
\pi(X_N)\,\mathrm dX_N
\propto \frac{z^N}{N!}\exp[-u(X_N)/T]
\prod_{m=1}^N\mathrm d\mathbf R_m\,\mathrm d\nu(\omega_m).
\label{eq:06-target}
\end{equation}
加速路径沿用相同的活度、分子内部参考和力场。下面的 $u$ 始终指输入指定的有限 Ewald 或离散 PME 能量；参考势 $u_0$ 只控制提议与筛选。

\subsection{固定框架的倒空间交叉项}

晶格向量为矩阵 $H$ 的行，$\mathbf r=\mathbf fH$，$V=\det H$，$\mathbf k_{\mathbf n}=2\pi\mathbf nH^{-\mathsf T}$。令 $C=e^2/(4\pi\epsilon_0 k_{\mathrm B})$，电荷 $q_i$ 以基本电荷为单位。对包含正负波矢的有限集合 $\mathcal K$，记
\begin{equation}
A_{\mathbf k}=C\frac{2\pi}{V}\frac{e^{-k^2/(4\alpha^2)}}{k^2},\qquad
S_H(\mathbf k)=\sum_{h\in H}q_h e^{\mathrm i\mathbf k\cdot\mathbf r_h},\qquad
S_G(\mathbf k)=\sum_{g\in G}q_g e^{\mathrm i\mathbf k\cdot\mathbf r_g}.
\label{eq:06-structure}
\end{equation}
减去空框架的倒空间能量后，精确分解为
\begin{align}
u_{\mathrm{rec}}^{\mathrm{rel}}
&=\sum_{\mathbf k\in\mathcal K} A_{\mathbf k}
\left[2\operatorname{Re}(S_H^*S_G)+|S_G|^2\right],\nonumber\\
u_{HG}^{\mathrm{rec}}
&=\sum_{g\in G}q_g\phi_H^{\mathrm{rec}}(\mathbf r_g),\qquad
\phi_H^{\mathrm{rec}}(\mathbf r)
=2\sum_{\mathbf k\in\mathcal K}A_{\mathbf k}
\operatorname{Re}\!\left[S_H(\mathbf k)^*e^{\mathrm i\mathbf k\cdot\mathbf r}\right].
\label{eq:06-host-split}
\end{align}
固定框架对任意位置的倒空间作用因此是一个标量场。框架 LJ 势按客体的 $(\sigma,\epsilon)$ 类型各存一份；实空间库仑与式~\eqref{eq:06-host-split} 合并为一个单位电荷场。类型 $t$ 在位置 $\mathbf r$ 的框架势为
\begin{equation}
v_t(\mathbf r)
=v_{t,H}^{\mathrm{LJ}}(\mathbf r)
+q_t\left[\phi_H^{\mathrm{real}}(\mathbf r)+\phi_H^{\mathrm{rec}}(\mathbf r)\right].
\label{eq:06-host-field}
\end{equation}
这种组织把相同 LJ 类型与不同部分电荷之间的冗余拆开，也使完整分子和单个 CBMC 生长位点共用同一批字段。

客体自能和净电荷背景仍由目标状态维护。三维 tin-foil 约定下，背景相对空框架的变化为
\begin{equation}
u_{\mathrm{self}}=-C\frac{\alpha}{\sqrt\pi}\sum_{g\in G}q_g^2,\qquad
u_{\mathrm{bg}}^{\mathrm{rel}}
=-C\frac{\pi}{2\alpha^2V}\left(2Q_HQ_G+Q_G^2\right).
\label{eq:06-self-background}
\end{equation}
式~\eqref{eq:06-host-field} 中没有这两项；它们通过第二阶段的完整能量校正进入接受率。

\section{一次 Fourier 变换生成框架场，空间缓存服务全部候选}
\label{sec:06-grid}

\subsection{有限 Fourier 和在网格节点上的精确抽样}

FFT 将周期 Fourier 展开转为空间场，是粒子网格长程算法的重要计算工具\cite{darden1993}。本版 direct 路径直接使用已缓存的 $S_H$ 生成场，不再为每个节点重新遍历框架电荷。设网格为 $M_x\times M_y\times M_z$，总点数 $M$，节点分数坐标为 $\mathbf f_{\mathbf j}=\mathbf j/\mathbf M$。把有限 Fourier 系数按各维模 $M_d$ 累加到频率数组：
\begin{equation}
\widehat\phi_{\mathbf n\bmod\mathbf M}
\mathrel{+}=2A_{\mathbf k_{\mathbf n}}S_H(\mathbf k_{\mathbf n})^*,\qquad
\phi_H^{\mathrm{rec}}(\mathbf f_{\mathbf j}H)
=M\,\operatorname{Re}\!\left[\operatorname{IFFT}(\widehat\phi)\right]_{\mathbf j}.
\label{eq:06-host-fft}
\end{equation}
这里 IFFT 包含 $1/M$ 归一化。频率采用累加而非覆盖，即使多个频率映射到同一网格索引，节点值仍等于原有限 Fourier 和。网格间距影响节点之间的插值质量，不会在建场时隐式删除原有 Ewald 模式。PME 路径则缓存固定框架的原生网格电势，并使用该 PME 模型的样条赋值规则生成相应的节点场。

LJ 与实空间库仑节点由现有原生邻域核分块计算；倒空间场只做一次 FFT。建场完成后，一个位点的普通查询只涉及周围 $4^3$ 个节点，生产阶段的框架评价成本与框架原子数脱离。

\subsection{周期三次插值与有限参考场}

周期网格使用分数坐标，天然适配三斜晶胞。每一维的局部坐标 $t\in[0,1)$ 对应四个 cardinal cubic 权重
\begin{align}
w_{-1}(t)&=-\tfrac12t+t^2-\tfrac12t^3,&
w_0(t)&=1-\tfrac52t^2+\tfrac32t^3,\nonumber\\
w_1(t)&=\tfrac12t+2t^2-\tfrac32t^3,&
w_2(t)&=-\tfrac12t^2+\tfrac12t^3.
\label{eq:06-cubic}
\end{align}
三维值为三个方向权重的张量积加权和。这是由邻近节点值构造的 Catmull--Rom 插值；能量和力网格在多孔材料模拟中的系统研究为网格分辨率与插值阶数的选择提供了背景\cite{ran2025}。

靠近框架原子核，LJ 排斥急剧增长，三次插值的负权重可能放大局部过冲。\code{HostGrid} 在预处理阶段按类型标记陡峭模板：普通区域使用式~\eqref{eq:06-cubic}，陡峭区域使用八节点、正权重的局部插值。正的 \code{host\_grid\_cap} 将节点参考值限制在有限区间，奇异节点取有限排斥值；记处理后的节点值为 $\bar v_a$，三线性权重为 $w_a\ge0$、$\sum_a w_a=1$。

默认陡峭模板插值 Boltzmann 权重：
\begin{align}
\widetilde v_B(\mathbf r;T)
&=-T\ln\sum_a w_a(\mathbf r)e^{-\bar v_a/T}\nonumber\\
&=v_{\min}-T\ln\sum_{a\in\mathcal J}w_a
e^{-(\bar v_a-v_{\min})/T},\quad
\mathcal J=\{a:w_a>0\},\quad v_{\min}=\min_{a\in\mathcal J}\bar v_a.
\label{eq:06-boltzmann-grid}
\end{align}
零权重角点不参与局部最小值，因而不会在节点或边界上引入无关的指数平移。由 Jensen 不等式，
\begin{equation}
\min_{a\in\mathcal J}\bar v_a
\le\widetilde v_B(\mathbf r;T)
\le\sum_a w_a(\mathbf r)\bar v_a.
\label{eq:06-jensen}
\end{equation}
能量线性插值会把一个高排斥角点按其几何权重加到整个网格单元；Boltzmann 插值让该角点的贡献指数衰减。如果一个低能角点主导局部和，其贡献约为 $\bar v_a-T\ln w_a$，从而保留粗网格单元内可进入区域的候选机会。\code{host\_grid\_steep=linear} 保留能量线性插值供对照。参考上限和插值均只定义 $u_0$；第二阶段由 $R$ 恢复目标势。上限设为零时，陡峭区域改用框架势的直接评价。

预计算的原始能量字段、奇异节点标志及物理参数指纹可以写入持久缓存。缓存键包含框架坐标与电荷、晶胞、客体 LJ 类型和电荷、截断与 Ewald 参数、PME/网格尺寸及参考上限；温度、储库活度与随机种子不参与原始字段定义。在这些势能参数固定时，同一文件可供不同温度、压力和随机种子的任务复用。式~\eqref{eq:06-boltzmann-grid} 在查询时使用当前温度，温度相关的参考局部能量与根位点提议表在各引擎中重新生成。对每次固定温度采样，$u_0(\cdot;T)$ 始终是确定的参考势。

\section{保留同一长程能量的模式与平移优化}
\label{sec:06-spectral}

对于实电荷，$S(-\mathbf k)=S(\mathbf k)^*$，正负波矢的 $A_{\mathbf k}$ 相同。令 $\mathcal K_+$ 从每对中保留一个波矢，并定义 $\widetilde A_{\mathbf k}=2A_{\mathbf k}$，则
\begin{equation}
\sum_{\mathbf k\in\mathcal K}A_{\mathbf k}|S_{\mathbf k}|^2
=\sum_{\mathbf k\in\mathcal K_+}\widetilde A_{\mathbf k}|S_{\mathbf k}|^2.
\label{eq:06-half-spectrum}
\end{equation}
默认半谱只改变存储和运算次数。节点场生成时向正负频率分别散射共轭系数，仍得到式~\eqref{eq:06-host-fft} 的同一场。若显式设置物理波数截止 $k_c$，模式集合再与 $|\mathbf k|\le k_c$ 相交；这一设置定义新的有限和，并进入配置指纹。

可选的离线建议工具进一步给出有限矩形集合内的截断界。设分子的绝对电荷和 $Q_*=\sum_{i\in m}|q_i|$，环境至多含 $N_{\max}$ 个同种分子，丢弃模式集合为 $\mathcal D$。由 $|S_B|\le N_{\max}Q_*$、$|\Delta S|\le Q_*$ 及三角不等式，完整分子的插入能量截断差满足
\begin{equation}
|\delta u_{\mathcal D}|
\le\sum_{\mathbf k\in\mathcal D} A_{\mathbf k}
\left[2\bigl(|S_H(\mathbf k)|+N_{\max}Q_*\bigr)Q_*+Q_*^2\right]
\equiv B_{\rm ins}.
\label{eq:06-tail-bound}
\end{equation}
同一环境中的旧、新完整分子分别满足此界，故平移、旋转和再生长的保守替换界为 $2B_{\rm ins}$。两个模式选择的实空间、自能、背景和分子内校正完全一致，在截断差中相消。\code{advise\_reciprocal.py} 按 $|\mathbf k|$ 壳层选择满足用户能差容限的最大丢弃集合，输出固定的 \code{ewald\_k\_cutoff} 建议；零表示保留原矩形集合。$N_{\max}$ 是离线分析前提，超过后需要重新求界；它不限制采样粒子数，也不用于生产过程中动态调节波矢。径向截断对应另一个具有所述能差界的有限 Hamiltonian，共轭半谱则保持原有限和不变。

每个接受分子保存自己的结构因子。新分子贡献和旧分子贡献之差 $\Delta S$ 给出
\begin{equation}
\Delta u_{\mathrm{rec}}
=\sum_{\mathbf k\in\mathcal K_+}\widetilde A_{\mathbf k}
\left[2\operatorname{Re}(S_{\mathbf k}^*\Delta S_{\mathbf k})+|\Delta S_{\mathbf k}|^2\right].
\label{eq:06-reciprocal-delta}
\end{equation}
对整体平移 $\mathbf d$，进一步利用
\begin{equation}
S_m'(\mathbf k)=e^{\mathrm i\mathbf k\cdot\mathbf d}S_m(\mathbf k),
\qquad Q_m'=Q_m,\qquad Q_{2,m}'=Q_{2,m}.
\label{eq:06-translation}
\end{equation}
因此平移无需再次遍历分子的全部位点。相位继续由三个基本方向的复数递推产生；位移取最终整体折回晶胞后的新旧锚点之差。旋转和一般再生长则构造新的完整分子贡献。

\section{参考势 CBMC：让每个 trial 保持局部计算}
\label{sec:06-reference-cbmc}

\subsection{分子内长程补全}

目标 Ewald 的分子内 primary pair 使用实空间项与 bare Coulomb 排除修正。若该对的库仑缩放系数为 $s_{ij}$，其进入短程分子内计算的部分为
\begin{equation}
u_{ij}^{\mathrm{intra}}
=Cq_iq_j\left[
\mathbf1_{r_{ij}<r_c}\frac{\operatorname{erfc}(\alpha r_{ij})}{r_{ij}}
-\frac{1-s_{ij}}{r_{ij}}\right].
\label{eq:06-intra-target}
\end{equation}
若仅删除候选阶段的倒空间计算，大的分子内排除项会单独留在参考权重中。为保留合理的局部能量尺度，参考势增加解析补全
\begin{equation}
u_{\mathrm{comp}}(m)
=C\sum_{i<j\in m}q_iq_j\frac{\operatorname{erf}(\alpha r_{ij})}{r_{ij}}.
\label{eq:06-intra-completion}
\end{equation}
当 $r_{ij}<r_c$，式~\eqref{eq:06-intra-target} 与补全之和恰为 $s_{ij}Cq_iq_j/r_{ij}$。完全排除对为零，1--4 对保留输入缩放，普通对恢复 primary bare Coulomb。$r_{ij}\ge r_c$ 时按同一代数分解定义参考项；$\operatorname{erf}(\alpha r)/r$ 在零点取有限极限。键、角、扭转和分子内 LJ 仍按原参数计算。

于是整个参考 Hamiltonian 可写为
\begin{equation}
u_0(X)
=\sum_{m}\sum_{i\in m}\widetilde v_{t_i}(\mathbf r_i)
+\sum_{m<n}u_{mn}^{\mathrm{short}}
+\sum_m\left[u_m^{\mathrm{intra}}+u_{\mathrm{comp}}(m)\right],
\label{eq:06-reference-energy}
\end{equation}
其中 $\widetilde v$ 是上述确定性框架参考场，$u^{\mathrm{short}}$ 保留客体间 LJ 与 Ewald 实空间库仑，$u^{\mathrm{intra}}$ 也包含键合能。所有项可随生长前缀逐段相加，每个相互作用在最后一个参与位点加入时计数。完整倒空间、自能、背景和网格残差集中在后续的 $R=u-u_0$ 中。分拆短程生长与长程校正的思路与既有 charged-CBMC 方法一致\cite{dubbeldam2013}。

\subsection{归一化权重、旧路径与根位点密度}

第 $i$ 阶段生成 $K_i$ 个候选，参考能量增量为 $\delta u_{0,ib}$，几何提议相对于内部参考测度的密度比为 $\rho_i$。使用
\begin{align}
\lambda_{ib}&=-\delta u_{0,ib}/T-\ln\rho_i(\mathbf r_{ib}),\nonumber\\
\ell_i&=\operatorname{LSE}_b(\lambda_{ib})-\ln K_i,
\qquad \ell=\sum_i\ell_i,\qquad W_0=e^\ell.
\label{eq:06-logmean}
\end{align}
候选选择概率正比于 $e^{\lambda_{ib}}$。\code{log\_weight} 使用减去最大值的 log-sum-exp；$K_i^{-1}$ 使改变 trial 数不会改变化学势。刚性分子把位置和取向作为一个阶段，柔性分子逐位点执行。删除和再生长的逆向权重由实际旧位点加 $K_i-1$ 个新候选计算，前缀来自实际旧分子。

可选的网格根位点提议为体素质量
\begin{equation}
p_v=\frac{1-\eta}{M}
+\eta\frac{\exp[-(\widetilde v_0(\mathbf r_v)-v_{\min})/T]}
{\sum_w\exp[-(\widetilde v_0(\mathbf r_w)-v_{\min})/T]},
\qquad 0\le\eta<1.
\label{eq:06-root-mixture}
\end{equation}
抽中体素后在其中均匀取点；alias 表把体素选择变为常数时间。均匀分量保留全部空间支持。由于每个体素体积为 $V/M$，根位点有 $\rho_0(\mathbf r)=Vq(\mathbf r)=Mp_{v(\mathbf r)}$；其余几何提议来自既定参考测度时 $\rho_i=1$。代码把 $T\ln\rho_0$ 加入候选权重，并在保存选中路径的参考物理能量时扣除它，因此提议密度不会混入残差 $R$。

\section{两阶段接受率与五类移动的详细平衡}
\label{sec:06-delayed-balance}

令 $\pi_0(X)\propto z^N e^{-u_0(X)/T}/N!$，定义 $R(X)=u(X)-u_0(X)$。对给定的前向候选集合及其对应逆向集合，CBMC 的参考对数比记为 $L_0$。目标比满足
\begin{equation}
L=L_0-\frac{R(Y)-R(X)}{T}.
\label{eq:06-factorization}
\end{equation}
Delayed acceptance 分别执行
\begin{equation}
a_1=\min(1,e^{L_0}),\qquad
a_2=\min\!\left(1,e^{-[R(Y)-R(X)]/T}\right),\qquad
a=a_1a_2.
\label{eq:06-two-stage}
\end{equation}
第一阶段拒绝时，不构造候选的完整长程贡献；通过后才评价式~\eqref{eq:06-factorization} 中的残差。该因子化接受规则遵循 delayed acceptance 的一般构造\cite{christen2005}。

为了直接对应实现，令 $B$ 为排除被操作分子后的同一环境，定义
\begin{equation}
e(m;B)=u(B\cup m)-u(B),\quad
e_0(m;B)=u_0(B\cup m)-u_0(B),\quad
r(m;B)=e(m;B)-e_0(m;B).
\label{eq:06-local-residual}
\end{equation}
记 $p_+$、$p_-$ 为插入和删除的移动选择概率，五类移动的两阶段对数比分别为：

\begin{center}\small
\begin{tabular}{@{}p{19mm}p{77mm}p{52mm}@{}}\toprule
移动 & 第一阶段 $L_0$ & 第二阶段 $L_R$\\\midrule
插入 & $\displaystyle\ln\frac{zV}{N+1}+\ell_{\rm new}+\ln\frac{p_-}{p_+}$ & $-r_{\rm new}/T$\\[7pt]
删除 & $\displaystyle\ln\frac{N}{zV}-\ell_{\rm old}+\ln\frac{p_+}{p_-}$ & $+r_{\rm old}/T$\\[7pt]
平移 & $-(e_{0,\rm new}-e_{0,\rm old})/T$ & $-(r_{\rm new}-r_{\rm old})/T$\\[4pt]
旋转 & $-(e_{0,\rm new}-e_{0,\rm old})/T$ & $-(r_{\rm new}-r_{\rm old})/T$\\[4pt]
再生长 & $\ell_{\rm new}-\ell_{\rm old}$ & $-(r_{\rm new}-r_{\rm old})/T$\\\bottomrule
\end{tabular}
\end{center}

平移和旋转提议关于新旧构型对称；插入与删除的 $N,V,z$ 和移动选择概率均显式保留。以候选集合和选中路径作为辅助变量，旧路径回溯给出一一对应的逆向构造。利用
\begin{equation}
\frac{\min(1,e^x)}{\min(1,e^{-x})}=e^x,
\end{equation}
两阶段前后向接受率之比为
\begin{equation}
\frac{a(X\to Y)}{a(Y\to X)}
=e^{L_0+L_R}
=\frac{\pi(Y)\,g(Y\to X)}{\pi(X)\,g(X\to Y)}.
\label{eq:06-detailed-balance}
\end{equation}
因此每类移动满足目标测度的详细平衡，按固定概率混合后仍保持式~\eqref{eq:06-target}。网格误差、节点上限及 root bias 可以改变候选质量和两阶段通过率；它们的作用由同一个 $R$ 校正。对于带净电荷的插入、删除和生长结果，第二阶段使用完整的 $Q$ 与 $Q_2$ 变化，式~\eqref{eq:06-self-background} 的交叉项随之保留。

\section{接受态缓存、恢复语义与计算成本}
\label{sec:06-execution}

接受态保存总结构因子、各分子贡献，以及每个分子的精确和参考局部能量。定义 $h(m)$ 为框架短程与分子内目标项之和，$h_0(m)$ 为网格框架项与分子内参考项之和。由于客体间 LJ 与 Ewald 实空间项在 $u$ 和 $u_0$ 中相同，残差直接化简为
\begin{equation}
r(m;B)=h(m)-h_0(m)+\Delta u_{\mathrm{rec+self+bg}}(m;B).
\label{eq:06-cancel-common}
\end{equation}
\code{selected\_residual} 先代数消去共有的客体间短程项，从而避免第二阶段重复访问客体邻域，也避免两个包含相同大项的完整能量相减。$h$ 和 $h_0$ 只依赖该分子和固定框架，其他分子移动时仍有效，因此旧分子的两项都可直接复用；倒空间交叉项使用当前环境。新分子贡献只构造一次，通过第二阶段后直接移动到接受态。

\begin{lstlisting}[basicstyle=\ttfamily\small]
choose a move; define the bath excluding the old molecule
build/retrace CBMC candidates using the reference field
if log(random()) >= min(0, L0): reject
evaluate selected residuals after cancelling common guest-short terms
form LR from the selected new/old residuals
if log(random()) >= min(0, LR): reject
commit coordinates, prepared charge contribution and both local energies
\end{lstlisting}

拒绝路径不改动接受态；框架场始终不可变。整体平移的相位递推会保留有限精度的运算历史，因而 checkpoint 同时保存总长程状态和逐分子 \code{ChargeContribution}，包括复结构因子或稀疏网格项、$Q$ 与 $Q_2$。恢复先建立同一模型，再覆盖为保存的数值缓存，并恢复随机数状态，使后续接受判据继承相同的浮点状态。

执行后端与评价粒度保持一致：\code{sampling\_mode=delayed} 的精确残差在原生 CPU 中计算，direct 路径采用半谱与相位更新；\code{sampling\_mode=exact} 保留对多个完整候选的 CUDA 倒空间批处理。前者把完整长程评价集中到少量选中构型，后者利用大候选批次的设备并行。

对 direct Ewald 路径，令 $A$ 为客体位点数，$B$ 为每阶段候选数，$K$ 为保留的倒空间模式数，$\bar n$ 为局部邻域规模。参考生长把主要候选工作变成网格插值、短程邻域及分子内局部计算；倒空间工作集中到第一阶段通过的少数完整构型。若第一、二阶段条件通过率分别为 $a_1,a_2$，平均单步成本可组织为
\begin{equation}
\mathbb E[C_{\rm step}]
\simeq C_{\rm reference}(A,B,\bar n)
+a_1C_{\rm selected\ residual}(A,K)
+a_1a_2C_{\rm commit}(K).
\label{eq:06-cost}
\end{equation}
普通 direct 候选的 $ABK$ 工作不再发生在每个 trial 上；平移的式~\eqref{eq:06-translation} 进一步将选中构型的结构因子更新降为 $K$ 级。网格建场成本由轨迹长度摊销，持久缓存则跨多个状态点复用。实际性能由参考场的区分能力、客体邻域、移动比例和两个通过率共同决定，因此运行诊断分别记录各移动的第一阶段次数、通过次数、精确评价次数与最终接受次数。

\begingroup\small
\begin{longtable}{@{}p{38mm}p{61mm}p{52mm}@{}}\toprule
数学对象 / 流程 & 实现位置 & 作用\\\midrule
固定框架场 & \code{HostGrid::build/value} & LJ 类型合并、单位电荷场、周期插值与陡峭模板。\\[3pt]
Fourier 节点场 & \code{EwaldState::host\_potential\_grid} & 原有限模式系数一次 FFT；PME 原离散场的独立缓存。\\[3pt]
共轭半谱与平移 & \code{EwaldState} 的模式行\newline\code{translated\_contribution} & 精确合并正负模式；用相位更新分子贡献。\\[3pt]
径向截断建议 & \code{tools06/advise\_reciprocal.py} & 按有限矩形尾和与占据上限给出固定参数建议。\\[3pt]
$u_0$ 生长增量 & \code{Fast::site\_energies}\newline\code{internal(...,true)} & 框架参考场、客体短程与分子内解析补全。\\[3pt]
$\ell$ 与旧路径 & \code{Fast::grow/log\_weight} & log-mean Rosenbluth、实际旧位点和提议密度修正。\\[3pt]
两阶段接受 & \code{Fast::step\_delayed} & 先参考比，后精确残差；覆盖五类移动。\\[3pt]
完整构型残差 & \code{Fast::selected\_residual} & 代数消去共有客体短程项；复用两种局部缓存与旧贡献。\\[3pt]
接受提交与恢复 & \code{commit\_insert/erase/replace}\newline\code{PoreMC-FastEwald-v3} & 缓存与坐标同步更新；保留逐分子浮点状态。\\[3pt]
跨状态点缓存 & \code{FastGCMC.\_grid\_cache\_key}\newline\code{HostGrid::load/save} & 物理输入指纹、完整性检查与一次写入后复用。\\\bottomrule
\end{longtable}
\endgroup

\clearpage
\section{独立数值与采样分布验证}

最终测试日志中的 335 项测试通过，另有 28 个参数子场景。新增核验同时比较能量恒等式和采样分布；后者直接检验参考势变化后，残差校正是否恢复目标分布。

\begin{center}\small
\begin{tabular}{@{}lr@{}}\toprule
验证对象 & 最大绝对误差\\\midrule
半谱总能 / 独立全谱 & $5.912\times10^{-12}$ K\\
增量能量 / 独立有限和 & $8.640\times10^{-12}$ K\\
框架 FFT 节点 / 直接求和 & $1.273\times10^{-11}$ K\\
平移相位 / 坐标重算结构因子 & $7.845\times10^{-15}$\\
Boltzmann 插值 / 独立加权配分和 & $1.990\times10^{-13}$ K\\
生产目标探针 / 独立 Ewald & $9.095\times10^{-12}$ K\\
\bottomrule\end{tabular}\end{center}

非均匀根位点检验把 85\% 的提议概率集中在 24 个 voxel 中的一个，同时保留全空间支持。理论计数均值和方差均为 2.7；16,000 个保存样本给出均值 2.71625、方差 2.73536。该算例检验根位点密度只进入 Rosenbluth 权重，而没有重复进入残差。

更强的检验故意把参考势截幅到 50 K，使其与带电单粒子的完整目标明显不同。独立空间积分和原生采样比较如下；误差条为分块标准误。

\begin{center}\small
\begin{tabular}{@{}lrrr@{}}\toprule
观测量 & 独立目标积分 & 未校正参考势 & 校正后采样\\\midrule
$\langle \cos(2\pi x/L_x)\rangle$ & -0.024514 & -0.001467 & $-0.013031\pm0.010045$\\
$\langle \sin(2\pi x/L_x)\rangle$ & -0.113245 & -0.007541 & $-0.109579\pm0.009006$\\
$\langle \cos(2\pi z/L_z)\rangle$ & 0.011903 & 0.003535 & $0.009988\pm0.009983$\\
\bottomrule\end{tabular}\end{center}

独立积分的网格加密最大变化为 $6.119\times10^{-5}$；本次实际残差范围为 -570.50 至 382.85 K。参考势本身不能复现目标分布，加入第二阶段后采样与目标积分一致。逐分子结构因子恢复、五类移动和缓存跨温度重用也已核验。

\section{完整循环的同机单线程基准}

\subsection{模型与计时约定}

CPU 为 AMD EPYC 9V74 80-Core Processor。两端使用相同编译优化，固定一个 OpenMP / BLAS 线程。对照源为打包的未修改 0.5 核心；0.6 所有主表均保持同一完整矩形有限 Ewald 目标，没有施加径向模式截断。

本轮没有用户实际 CIF 与完整 p-xylene 力场输入，因此创建并公开了计算规模匹配的合成算例：48 \AA{} 立方盒、4,416 个框架点、12 个客体位点、10 个 CBMC trial、10 个初始分子、$k_{\max}=13$、$T=350$ K、$r_c=10$ \AA{}、$\alpha=0.28$ \AA$^{-1}$。五种移动等概率，平移上限 0.65 \AA{}、转角上限 0.35 rad；LJ 混合采用 Lorentz--Berthelot 与截断移位，1--4 LJ / Coulomb 缩放为 0.5 / 0.75。

框架为形成周期狭缝的八个轻微起伏平面，位点 $\sigma=2.9$ \AA、$\epsilon=4$ K；电荷由固定种子生成并减去均值。刚性算例由六个环 C、四个 H 和两个甲基伪位点组成，采用公开在脚本中的合成参数；柔性算例为十二位点链，$\sigma=2.2$ \AA、$\epsilon=6$ K，电荷交替 $\pm0.03e$，键长 1.3 \AA、键角 $112^\circ$、三重扭转系数 15 K。两类活度分别为 $8/V$ 和 $12/V$。这些参数定义数值和性能任务，不是对某种材料的力场标定。

刚性每组预热 500 次、生产 3000 次；柔性分别为 200 与 1500 次。各三种子、独立进程并交替测试顺序。生产吞吐不含初始化和末态参考重算，关闭逐帧总能记录；接受次数与计数轨迹同时保留。

\subsection{十二位点刚性客体}

\begin{center}\small
\begin{tabular}{@{}lrrrr@{}}\toprule
路径 & 尝试/s & 接受/s & 配对加速 & 加速范围\\\midrule
0.5 全谱逐候选 & 316.10 & 197.14 & 1.00 & 1.00--1.00\\
0.6 半谱逐候选 & 523.38 & 326.42 & 1.67 & 1.66--1.70\\
0.6 参考网格 $32^3$ & 4,933.15 & 2,428.76 & 16.71 & 15.53--17.14\\
$32^3$ + alias 根位点 & 4,972.70 & 2,405.82 & 16.41 & 15.73--16.79\\
0.6 参考网格 $64^3$ & 4,499.63 & 2,507.79 & 14.23 & 11.27--16.23\\
\bottomrule\end{tabular}\end{center}

\subsection{十二位点柔性客体}

\begin{center}\small
\begin{tabular}{@{}lrrrr@{}}\toprule
路径 & 尝试/s & 接受/s & 配对加速 & 加速范围\\\midrule
0.5 全谱逐候选 & 101.32 & 63.83 & 1.00 & 1.00--1.00\\
0.6 半谱逐候选 & 186.20 & 116.90 & 1.85 & 1.83--1.92\\
0.6 参考网格 $32^3$ & 4,683.53 & 2,740.62 & 46.22 & 43.00--49.43\\
\bottomrule\end{tabular}\end{center}

吞吐两列分别为三次测量的中位数；加速列为每次配对比值的中位数，因此不必等于两列中位数之比。接受吞吐避免把纯拒绝路径的高尝试数误认成构型探索效率；构型相关性另在长链部分检验。

全部配对任务的末态完整重算最大差为 $9.286\times10^{-9}$ K。不同接受算法使用不同数量的随机数，轨迹无需逐步相同；物理输入哈希、初始能量、末态重算和分布共同定义比较依据。

\subsection{缓存初始化与生产工作量}

\begin{center}\small
\begin{tabular}{@{}lrrr@{}}\toprule
网格 & 首次构建/s & 读取复用/s & 磁盘字段/MiB\\\midrule
$32^3$ & 0.672 & 0.115 & 1.13\\
$64^3$ & 4.689 & 0.282 & 9.00\\
\bottomrule\end{tabular}\end{center}

这里的初始化包括势函数和原生引擎构建，单独记录，不加入生产吞吐。缓存保存原始字段，重新选择温度、活度和根位点偏置时可以继续读取。

\begin{center}\small
\begin{tabular}{@{}lrrr@{}}\toprule
路径 & 第一阶段通过/\% & 第二阶段通过/\% & 精确分子评价/尝试\\\midrule
0.6 参考网格 $32^3$ & 67.95 & 72.50 & 1.109\\
$32^3$ + alias 根位点 & 67.93 & 70.50 & 1.105\\
0.6 参考网格 $64^3$ & 66.69 & 89.07 & 1.089\\
\bottomrule\end{tabular}\end{center}

阶段统计包含预热与生产；空体系无法执行的删除等移动不进入第一阶段分母。一次替换包含旧、新两个残差，旧分子结构因子与两种局域能直接读取缓存。

\section{粗网格的负结果与混合检验}

初始三线性参考把非常高的排斥节点能量混入相邻可达体积。下表给出相同刚性任务三条短链的平均计数，保留失败的优化路径以解释最终设计。Boltzmann 插值恢复了进入近壁区域的机会，但参数选择仍以更长轨迹和相关性判断。

\begin{center}\small
\begin{tabular}{@{}lr@{}}\toprule
路径 & 三种子的平均粒子数\\\midrule
精确 0.5 参考 & 13.79, 11.97, 14.24\\
线性能量插值 $32^3$ & 9.45, 8.86, 8.63\\
Boltzmann $32^3$ & 11.44, 13.36, 10.58\\
Boltzmann $64^3$ & 11.10, 13.70, 13.30\\
\bottomrule\end{tabular}\end{center}

长链对比每种路径运行三个独立种子，每条预热 5,000 次、生产 30,000 次，每 10 次保存计数。下表的误差为三个种子均值的标准误；ESS 按各链的初始正序列自相关估计后相加，仅针对粒子数这一观测量。

\begin{center}\small
\begin{tabular}{@{}lrrr@{}}\toprule
路径 & $\langle N\rangle$ & 种子均值 SE & $N$ 的 ESS 合计\\\midrule
0.5 全谱逐候选 & 13.936 & 0.395 & 199.6\\
0.6 参考网格 $32^3$ & 13.799 & 0.708 & 119.4\\
0.6 参考网格 $64^3$ & 13.242 & 0.243 & 220.8\\
$32^3$ + alias 根位点 & 14.451 & 0.356 & 129.8\\
\bottomrule\end{tabular}\end{center}

相同 90,000 次生产尝试下，$64^3$ 路径的计数 ESS 与精确对照相近，且优于本例的 $32^3$ 路径；结合完整循环吞吐，本报告建议从 $64^3$ Boltzmann 参考网格开始。$32^3$ 与 alias 参数仍可按目标体系实测选择。独立链均值仍有散布，表中没有把这组诊断作为精确等温线的最终物性。长链为缩短验证时间并行运行，其耗时不用于性能倍率或 ESS/s。

\section{显式模式截断的精度与成本空间}

除保持原有限和的半谱优化外，本版提供离线误差界建议工具。给定最大环境载量与完整分子的电荷绝对值和，可以界定省略倒空间壳层造成的能量增量差。下面使用同一刚性合成模型、$N_{\max}=64$；保守替换界覆盖完整分子的插入、删除与替换。主性能表没有采用这些额外截断。

\begin{center}\small
\begin{tabular}{@{}lrrrr@{}}\toprule
容限/K & $k_c$/\AA$^{-1}$ & 保留半谱 & 替换上界/K & 实测最大差/K\\\midrule
0.01 & 2.306566 & 8851 & 0.009744 & $1.162\times10^{-5}$\\
0.1 & 2.168751 & 8193 & 0.097322 & $7.232\times10^{-5}$\\
1 & 2.038423 & 7289 & 0.973553 & $6.917\times10^{-4}$\\
\bottomrule\end{tabular}\end{center}

每个建议另以环境占据 0、10、64 的 12 组新旧刚体构型检查实际省略差。解析界是给定载量范围内的约束，样本最大差只是这组构型的测量。非零径向截断定义另一个有限 Ewald Hamiltonian；没有把其模式数缩减当作本轮已测得的额外加速。

\begin{lstlisting}
python tools06/advise_reciprocal.py your_hpc_input.json \
    --max-molecules 64 --tolerance-K 0.1 \
    --output reciprocal_advice.json
\end{lstlisting}

该工具只输出建议，不修改输入文件。载量超出分析范围时应重新评估，生产采样中不动态切换目标模式集合。


\section{运行与真实 HPC 输入的复测}
\subsection{构建和直接运行}
\begin{lstlisting}
python -m pip install -e ".[test]"
python native/build.py
python -m poremc run examples06/small_delayed_ewald.json -o run06
\end{lstlisting}
原有输入保持力场、活度与 Ewald 参数，在 \code{sampling} 中加入：
\begin{lstlisting}
"driver": "cpp",
"sampling_mode": "delayed",
"host_grid_mesh": [64, 64, 64],
"host_grid_steep": "boltzmann",
"host_grid_cap": 100000.0,
"host_grid_cache": "cache/framework-guest-fields.bin",
"grid_root_bias": 0.0,
"ewald_half_spectrum": true,
"ewald_k_cutoff": 0.0,
"reciprocal_method": "direct",
"reciprocal_backend": "cpu"
\end{lstlisting}
\code{host\_grid\_cache} 相对输入文件所在目录解析。复用同一框架与力场的等温线或温度扫描时，原始网格字段可以直接读取；温度相关的 Boltzmann 插值和根位点概率在当前实例中计算。

\code{grid\_root\_bias=0.95} 启用混合 alias 提议，其中 5\% 保留均匀空间支持；它是可选优化。把 \code{host\_grid\_steep} 设置为 \code{'linear'}，可复现初始三线性参考的粗网格混合反例。

把 \code{sampling\_mode} 设置为 \code{'exact'}，使用逐候选精确能量评价，并保留原 CUDA 批处理入口。新的 delayed 校正当前在原生 CPU 执行；\code{auto} 选择 CPU，显式 \code{cuda} 请求需使用 exact 模式。

\subsection{在用户实际构型上比较版本}
以下命令读取完整 PoreMC JSON，包括 CIF、分子力场和初始分子，自动编译打包的原始 0.5 对照，按交替顺序测量多个版本：
\begin{lstlisting}
python tools06/benchmark_input.py your_hpc_input.json \
    --steps 3000000 --burnin 10000 --repeats 3 \
    --modes baseline05,native06,delayed32,biased32,delayed64 \
    --output hpc_comparison_06.json
\end{lstlisting}
该命令的 \code{steps} 表示尝试移动数。所有变体固定单线程，并采用输入的原矩形 direct Ewald 模式集合；记录初始化、生产循环、接受次数、载量、实际后端和末态重算差。程序不把本机倍率乘到用户原 HPC 数值上作为新的实测，也不将 PoreMC 的尝试数自动写成 RASPA cycles。

\subsection{诊断和继续优化的接口}
\code{simulation.acceleration\_profile} 返回各移动的第一阶段通过数、精确残差评价数、第二阶段接受数，以及网格命中和插值类型。\code{long\_range\_profile} 给出实际存储模式、完整模式、FFT、结构因子构造与 GPU batch 次数。

\code{trial\_molecule\_energies} 与 \code{trial\_site\_energies} 始终返回目标能量；对应 \code{trial\_surrogate\_...} 接口返回参考势，便于研究新的网格、学习势和分子提议。checkpoint v3 保留逐分子结构因子及长程累计状态，使平移相位递推后的轨迹可以逐值恢复。

原生范围仍为单组分刚性 / 树状柔性分子和三维周期静电。现有 0.3/0.4 电解质、恒电位电极、反应与电流接口继续保留；本轮变化集中在 \code{FastGCMC} 的固定框架采样路径。

\section{复现材料}
\begingroup\small
\begin{longtable}{@{}p{0.45\linewidth}p{0.49\linewidth}@{}}\toprule
路径 & 内容\\\midrule
\code{native/host\_grid.*} & 字段建表、插值、Boltzmann 参考与持久缓存。\\
\code{native/ewald\_state.*} & 半谱、host FFT 场、平移相位和增量静电。\\
\code{native/fast\_mc.cpp} & 参考势 CBMC、两阶段接受、残差相消与缓存提交。\\
\code{tools06/benchmark\_hostgrid.py} & 同规模合成算例的完整循环配对计时。\\
\code{tools06/benchmark\_input.py} & 用户真实输入的同物理模型版本比较。\\
\code{tools06/benchmark\_mixing.py} & 多种子长链载量与相关性检验。\\
\code{tools06/audit\_delayed.py} & 独立物理参考、非均匀提议与分布核验。\\
\code{tools06/advise\_reciprocal.py} & 给定载量范围和能量界的显式模式建议。\\
\code{baseline\_05/} & 未修改的 0.5 核心代码与原交付来源。\\
\code{evidence06/} & 实测 JSON、测试日志、负结果和来源说明。\\\bottomrule
\end{longtable}
\endgroup
\begin{lstlisting}
python -m pytest -q tests tests_legacy tests03 tests04 tests06
python tools06/audit_delayed.py
python tools06/audit_host_grid.py
python report_06/generate_results.py --require-complete
python report_06/build_report.py
\end{lstlisting}
网格二进制缓存可由模型重新构建，不作为源码依赖打包。源文件与原始证据的哈希保存在发布清单中。

\begingroup\small
\begin{thebibliography}{9}
\setlength{\itemsep}{3pt}\setlength{\parskip}{0pt}
\bibitem{christen2005} J. A. Christen, C. Fox.
Markov chain Monte Carlo Using an Approximation.
\textit{J. Comput. Graph. Stat.} \textbf{14}(4), 795--810 (2005).
\href{https://doi.org/10.1198/106186005X76983}{doi:10.1198/106186005X76983}.
\bibitem{dubbeldam2013} D. Dubbeldam, A. Torres-Knoop, K. S. Walton.
On the inner workings of Monte Carlo codes.
\textit{Mol. Simul.} \textbf{39}(14--15), 1253--1292 (2013).
\href{https://doi.org/10.1080/08927022.2013.819102}{doi:10.1080/08927022.2013.819102}.
\bibitem{ran2025} Y. A. Ran et al.
Implementation of energy and force grids in molecular simulations of porous materials.
\textit{Mol. Phys.} \textbf{123}(21--22), e2542486 (2025).
\href{https://doi.org/10.1080/00268976.2025.2542486}{doi:10.1080/00268976.2025.2542486}.
\bibitem{darden1993} T. Darden, D. York, L. Pedersen.
Particle mesh Ewald: An $N\log(N)$ method for Ewald sums in large systems.
\textit{J. Chem. Phys.} \textbf{98}, 10089--10092 (1993).
\href{https://doi.org/10.1063/1.464397}{doi:10.1063/1.464397}.
\end{thebibliography}
\endgroup
\end{document}
